\section{Field and fluence dependences of laser-induced multiple spin-wave dynamics in $\mathrm{Co_2FeAl_{0.5}Si_{0.5}}$ films \href{https://doi.org/10.1063/5.0006321}{[\textit{J. Appl. Phys., 127 233903 (2020)}]}}


\subsection*{\underline{Abstract}}
Field and fluence-dependent spin-wave dynamics of a full-Heusler $\mathrm{Co_2FeAl_{0.5}Si_{0.5}}$ film studied by using TR-MOKE. Volumetric magnetostatic spin-wave (VMSW) and perpendicular standing spin-wave (PSSW) modes are excited in the films with thicknesses of 60 and 100 nm. Only Kittel mode is observed in the films of thicknesses of 150 and 200 nm. The amplitude of all spin-waves increased with increasing field and fluence and frequency slightly decrease with increasing fluence as expected. Effective damping and lifetimes are both modulated by both external field and excitation fluence. This changes are attributed to the field-dependent group velocity and magnetic inhomogeneity. The dominant mechanism of the electron-phonon scattering enhancement might be the reason for the increment of the effective damping with increasing fluence.

\subsection*{\underline{I. Introduction}}
\begin{enumerate}
    \item Cobalt based heusler/alloys exhibit strong application potentials for magnonic technologies.
    \begin{itemize}
        \item High spin polarization.
        \item Low intrinsic magnetic damping.
        \item Possibility of exciting multiple spin-wave modes (SWMs).
    \end{itemize}
    \item Different spin waves including uniform spin wave (Kittel mode), magnetostatic surface spin-wave (Damon-Eshback (DE) mode), volume magnetostatic spin wave (VMSW), and perpendicular standing spin wave (PSSW) were coherently excited and identified in different cobalt-based Heusler alloy films.
    \item The field-and fluence-dependent spin wave dynamics in full-Heusler $\mathrm{Co_2FeAl_{0.5}Si_{0.5}}$ films with different thicknesses by using TR-MOKE experimental setup with an out-of-plane external field applied.
    \item The VMSW and PSSW modes were excited in two thinner samples, while in the two thicker samples only the Kittel mode is excited.
\end{enumerate}

\subsection*{\underline{II. Experiments}}
\begin{enumerate}
    \item Magnetic sputtering was used to deposit the CoFeAlSi sample on the glass substrate.
    \item The thickness of the sample is 60, 100, 150, and 200 nm.
    \item All the made samples have in-plane magnetized feature that can be attributed to the demagnetizing field.
    \item Pump-probe was used to understand the spin-wave dynamics.
    \item A variable magnetic field generated by an electromagnet is applied along the direction of $\sim 10^\circ$ with respect to the normal of the sample plane.
    \item A nearly perpendicular orientation of the external field is beneficial to create a large precession angle after laser pumping and excite additional SWMs for in-plane magnetized samples [\textcolor{red}{ref. 21}].
    \item All measurements are performed at room temperature $[\mathrm{\sim300 K}]$.
\end{enumerate}
\subsection*{\underline{III. Results and Discussions}}
\subsubsection*{A. Spin-wave dynamics and dispersion analysis}
\begin{enumerate}
    \item The precession of magnetization is launched by the torque exerted on it as the femtosecond pump pulse changes the orientation of the effective field transiently [\textcolor{red}{ref. 19}].
    \item Oscillations with large amplitude with respect to demagnetization indicate the efficient excitation of spin waves.
    \item The Fast Fourier Transform (FFT) spectra of the oscillatory components of the dynamics are obtained and it clearly showed the existence of the spin-waves for 60 nm and 100 nm thickness.
    \item The FFT spectra clearly indicate two different SWMs in the 60 and 100 nm samples but only one SWM in the 150 and 200 nm samples is excited.
    \item The dispersion analysis was also performed to identify the excited SWMs.
    \item Competition between the out-of-plane external field and the demagnetization field results in the slant orientation of the equilibrium effective field.
    \item For the SWM (Kittel mode), the dispersion equation is written as;
    \begin{gather}
        \mathrm{\omega^2_{Kittel}=\omega_H(\omega_H+\omega_M\sin^2\theta)}
    \end{gather}
    where $\mathrm{\omega_H=\gamma H\sin \phi/\sin\theta}$ and $\mathrm{\omega_M=4\pi\gamma M_s}$, in which $\theta$ and $\phi$ are the orientation angles of equilibrium magnetization and the external field with respect to the normal of the film plane, respectively.
    \item On the other hand, for the exchange-interaction dominated PSSW mode, the dispersion equation is given as [\textcolor{red}{ref. 22}];
    \begin{gather}
        \mathrm{\omega^2_{PSSW}=(\omega_H+\omega An^2)(\omega_H+\omega_M\sin^2\theta+\omega_A n^2)}
    \end{gather}
    where $\mathrm{\omega_A=2\gamma A_{ex}(\pi/d)^2/M_s}$, in which $A_{ex}$ is the exchange constant and $n$ denotes the order of the PSSW mode.
    \item Overall, from the analysis it can be obtained that the VMSW and PSSW can be assigned to the LF and HF modes, respectively, in the cases of 60 and 100 nm, while the Kittel mode can be assigned to the sole SWM in the cases of 150 and 200 nm.
\end{enumerate}
\subsubsection*{B. Field dependence of lifetimes, damping factors, and amplitudes}
\begin{enumerate}
    \item We were able to reveal the energy dissipation features of spin waves with the help of lifetime and damping calculations.
    \item In the case of 60 and 100 nm, the lifetimes of the VMSW mode mainly decrease with increasing H.
    \item The lifetimes of the PSSW mode are apparently lower than those of the VMSW mode (which was around 862 ps for 60 nm in the H-field of 0.8 kOe).
    \item In the case of 150 and 200 nm, the lifetimes of the Kittel mode increase with increasing H, showing an opposite trend with respect to the field dependence of the VMSW mode in the cases of 60 nm and 100 nm.
    \item As expected, the effective damping factors of the Kittel mode in the cases of 150 and 200 nm present an opposite trend with respect to those of the VMSW mode in the case of 60 and 100 nm.
    \item The effective damping of the VMSW mode generally increases with H first and then decreases slightly in the high field range, showing the characteristic of a back VMSW mode excited in an in-plane magnetized film under a perpendicular external field.
    \item This field-dependence of the effective damping can be attributed to the energy propagation away from the probe area, which is mainly described by the group velocity strongly depending on the external field [\textcolor{red}{ref. 15}].
    \item The main extrinsic contribution to such field-dependent damping is magnetic inhomogeneity.
    \item This can be caused by the non-uniform morphology, defects, or domain formation at low field which could result in a distribution of magnetic property within the probe area.
    \item The contribution of magnetic inhomogeneity to the effective field would be reduced in a high field range, leading to a lower effective damping [\textcolor{red}{19, 26}].
    \item The differences in the effective damping parameter for the 60 nm and 100 nm in the case of VMSW and 150 nm and 200 nm for the PSSW in the external field of $H=8$ kOe and a pump fluence of 22.5 $\mathrm{mJ/cm^2}$, is not much and it can be implied that the influence of magnetic inhomogeneity can be ignored in the cases of 60 nm and 100 nm.
    \item Additionally, the effective damping factors of the PSSW mode approximately remain stable with the increase in $H$.
    \item Moreover, as one can note that the pump-fluence dependences of effective damping for the VMSW and PSSW modes also surprisingly present an opposite trend with respect to that for the Kittel mode.
    \item The extracted amplitudes of the spin waves of the 60 and 200 nm as the functions of $H$ and with different pump fluence.
    \item As expected, the amplitudes of all the three modes increase with $H$, which can be attributed to the larger out-of-plane magnetization component under the higher out-of-plane external field.
\end{enumerate}
\subsubsection*{C. Fluence dependence of dynamics parameters}
\begin{enumerate}
    \item The trends in the fluence dependences of the spin-wave dynamics, the parameters obtained under $H$ of 8 kOe in the cases of 60 and 200 nm have the same variation trends as those for 60 and 200 nm.
    \item As expected, the amplitudes of all three modes increases with increasing pump fluence, which originates from the stronger laser modulation to magnetization and the effective field.
    \item And in the case of frequencies, there is a slight decrement as the pump fluence is increased, which is in accordance with the fluence dependence of the frequency for the Kittel mode in $\mathrm{Co_2MnSi}$ film [\textcolor{red}{ref. 13}].
    \item This decrement in the frequency is attributed to the rise in the equilibrium temperature, as the main contribution from the demagnetizing field, $\omega_M \sin^2\theta$ is decreased and cause this effect.
    \item Also, as the pump-fluence is increased, the effective damping increases. This is in agreement also as per the DE and PSSW modes obtained in the $\mathrm{Co_2FeAl}$ film and can be similarly explained by the enhancement of electron-phonon scattering.
    \item As the fluence is increased, the energy dissipation from the spin system to lattice is increased and hence, this accelerates the decay of spin waves. This was demonstrated by the measurement of coherent acoustic phonon [\textcolor{red}{ref. 9,12}].
    \item As for Kittel mode, lifetime increases but effective damping decreases with increasing pump fluence. This is in contrast with the VMSW and PSSW modes but agree with the pump-fluence-dependent damping feature of Kittel spin precession observed in some other ferromagnetic films such as FePt [\textcolor{red}{19}].
    \item The magnetic inhomogeneity is the dominant extrinsic damping origin.
    \item Generally, the influence of magnetic inhomogeneity on spin precessions can be evaluated based on two aspects, the dispersion of internal field (effective anisotropy field) and the contribution of the inhomogeneous internal field to the precession lifetime [\textcolor{red}{29}].
    \item Hence, it is easily understood that a higher pump fluence can suppress the frequency broadening and lead to lower effective damping.
    \item The PSSW mode also cannot be observed in these two thicker samples. This may be because the frequencies of PSSW for the thicknesses of 150 and 200 nm are already close to those of the Kittel mode, and thus, the laser pulses cannot simultaneously excite such tow SWMs.
\end{enumerate}